\section{Differenzialrechnung von Funktionen einer
Veränderlichen}

\begin{ejercicio}[Ableitung trigonometrischer Funktionen mit Produktregel]
    Differenzieren Sie die folgenden Funktionen nach der Variablen $x$:
    \begin{enumerate}
        \item $f(x) = 3x^2 \cdot \cos(x)$
        \item $f(x) = \sin(x) \cdot \cos(x)$
    \end{enumerate}
\end{ejercicio}

\begin{ejercicio}[Quotientenregel]
    Differenzieren Sie die folgenden Funktionen nach der Variablen $x$ und bestimmen Sie den größtmöglichen Definitionsbereich in $\mathbb{R}$:
    \begin{enumerate}
        \item $f(x) = 2 \cdot \frac{\ln(x)}{x}$
        \item $f(x) = 2 \cdot \frac{\cos(x)}{3x^2}$
    \end{enumerate}
\end{ejercicio}

\begin{ejercicio}[Kettenregel]
    Differenzieren Sie nach der Variablen $x$:
    \begin{enumerate}
        \item $f(x) = 3 (2x + 1)^4$
        \item $f(x) = e^x$
        \item $f(x) = e^{2x}$
        \item $f(x) = e^{2x^3}$
    \end{enumerate}
\end{ejercicio}

\begin{ejercicio}[Gemischte Aufgaben zur Differenzialrechnung]
    Differenzieren Sie nach der Variablen $x$ und geben Sie ggf. Einschränkungen des Definitionsbereichs an, falls die Differenzierbarkeit nicht auf ganz $\mathbb{R}$ gegeben ist.
    \begin{enumerate}
        \item $f(x) = x^{x^2}$
        \item $f(x) = (x + 1)^x$
        \item $f(x) = x^{(x+1)}$
        \item $f(x) = (x + 1)^{(x+1)}$
        \item $f(x) = \left(2x^2+1\right)^{\sin(x)}$
        \item $f(x) = \frac{e^{-\cos x}}{\sin(x)}$
    \end{enumerate}
\end{ejercicio}

\begin{ejercicio}[Differenzieren in 1D - Test]
    Differenzieren Sie nach der Variablen $x$. Welche Bedingungen muss der jeweilige Definitionsbereich erfüllen?
    \begin{enumerate}
        \item $f(x) = \left(2x^2 - \frac{4}{x} + 3\right)^4$
        \item $f(x) = e^{2x \cos x}$
        \item $f(x) = \ln\left(\frac{1}{x^2}\right)$
        \item $f(x) = \sqrt{\frac{x^2 - a^2}{a^2 + x^2}}$
    \end{enumerate}
\end{ejercicio}

\begin{ejercicio}[Kurvendiskussion]
    Diskutieren Sie den Verlauf der Funktion
    \[
        f(x) = \frac{(x - 2)^2}{x + 2}.
    \]
    Machen Sie Aussagen zu Definitionsbereich, Symmetrie, Nullstellen, Polen, Ableitungen (bis zur 3. Ordnung), relative Extremwerte, Wende- und Sattelpunkte und dem Verhalten ``im Unendlichen''. Fertigen Sie am Schluss eine saubere Skizze des Kurvenverlaufs an.
\end{ejercicio}